\section{Resilience and Climate Outcomes}

\subsection{Compound Exposure Hardening: Closure Cost Reduction}

The Track D incident economics analysis estimates current annual closure costs at \$6.2 billion per year across all T1 and T2 corridors, with the top seven compound exposure corridors accounting for \$3.9 billion of that total \citep{ROUTE_D2}. The compound exposure hardening program (Component 3, \$18B capital) is designed to reduce that \$3.9 billion to approximately zero through targeted physical interventions: flood elevation of roadway to 2050 projected flood stage plus 18 inches freeboard, fire-resistant surface treatments on segments with high wildfire interface risk, hurricane wind load upgrades on elevated structures within the Gulf Coast storm surge zone, and secondary drainage channels on segments with recurring flash flood closures.

The \$3.9B/yr closure cost reduction is the primary annual benefit stream for Component 3. At 7\% discount over 30 years, this produces a \$37.8B NPV against the \$18B capital cost --- a 3.1:1 B/C ratio \citep{ROUTE_B3,ROUTE_D2}. This is the highest per-dollar return in the portfolio because hardening addresses both dimensions of the compound exposure condition simultaneously: a corridor hardened against flood is also made more structurally resilient against the secondary failure modes (undermined foundations, scoured abutments) that extend closures from days to weeks.

\subsection{Diamond Interchange Upgrades: Eliminating k=1 Conditions}

Nine of the 15 T1/T1 intersections currently have k=1 connectivity: no viable alternate route exists if the primary interchange is closed \citep{ROUTE_B4}. The diamond interchange upgrade program (Component 2, \$4.5B) eliminates the k=1 condition at all 15 nodes by providing grade separation, signal-free freight movement, and alternate ramp configurations that remain operable when one ramp set is closed for maintenance or incident management.

The Phase 1 diamond upgrade portfolio alone (six highest-priority nodes, approximately \$1.6B) prevents an estimated \$6.5B in incident costs during the 2025--2035 decade. This figure is derived from the Track B incident cost model, which estimates the expected annual cost of k=1 closure events (probability $\times$ consequence) at each node using historical incident frequency data from FHWA's Traffic Monitoring Guide and freight value from the BTS Commodity Flow Survey \citep{FHWA_freight2023}.

\subsection{2050 Gulf Coast Climate Resilience}

The Gulf Coast I-10 corridor --- from Baton Rouge, LA to Pensacola, FL --- is the highest-priority compound exposure investment in the portfolio. Its projected D1 score in 2050 is 9.1, driven by three compounding factors: (1) sea-level rise of 1.4--2.2 feet (median NOAA projection for the Gulf Coast by 2050 \citep{NOAA_SLR2022}), (2) increasing Category 4 and 5 hurricane probability with post-storm surge extending inland to the I-10 alignment in multiple Louisiana and Mississippi parishes, and (3) the ongoing subsidence of coastal Louisiana at 0.4--0.9 inches per year, which reduces the effective flood protection provided by existing levee systems.

Without hardening intervention, Track D projects that private freight carriers --- particularly less-than-truckload carriers operating on fixed schedule commitments --- will begin systematically rerouting around Gulf Coast I-10 before 2045, diverting to I-20 and I-40 as the primary east-west southern trunk. This rerouting effectively retires Gulf Coast I-10 as a T1 artery, adding 140--200 miles to the southeastern freight path and increasing NY--Miami transit times by 4--6 hours on average \citep{ROUTE_D1}. The \$3.6 billion Gulf Coast I-10 hardening investment in Phase 1 prevents this outcome, maintaining T1 functionality on the corridor through 2050 and beyond.

\subsection{Missing Link Redundancy Value}

The five proposed missing link corridors contribute resilience value beyond their primary freight and connectivity benefits:

\textbf{I-70W/US-50 (Donner redundancy)}: The Track D analysis estimates that Donner Pass closures currently cost \$2.4B/yr in freight disruption costs, considering both the direct rerouting cost and the reliability premium charged by carriers for transcontinental shipments with elevated Donner risk \citep{ROUTE_D2}. With I-70W/US-50 upgraded to interstate standard, the Donner closure consequence drops to approximately \$0.5B/yr, because drivers can legally use the interstate-standard US-50 alternative without the weight and height restrictions that currently limit this route for Class 8 freight. For comparison, the alternative proposal --- a Donner tunnel at approximately \$4 billion --- reduces the closure cost to \$0.8B/yr by shortening the high-risk mountain segment, but does not provide the systemic redundancy of a full alternate corridor.

\textbf{I-69 (Gulf--Midwest path diversity)}: I-69 completion reduces the I-35 concentration of Gulf--Midwest freight, decreasing the consequence of any I-35 closure event by approximately 25\% (proportional to the traffic diversion estimated in Track C's max-flow analysis \citep{ROUTE_C2}).

\subsection{Network Reliability: Variance Reduction Summary}

The combined effect of managed freight lanes (PTI reduction), diamond interchange upgrades (k=1 elimination), and compound exposure hardening (closure probability reduction) produces the reliability variance reduction described in Section~4. The 57\% reduction in PTI standard deviation across T1 corridors (from 0.42 to 0.18) translates to an estimated 40\% reduction in national freight reliability variance, measured as the coefficient of variation of shipper-reported transit times on the ten highest-volume O-D pairs \citep{ROUTE_C1}.

This variance reduction is the I2.0 program's most significant systemic benefit. The \$8.2B annual reliability cost identified in Track C is not primarily caused by mean congestion delay; it is driven by the unpredictability of delay. Shippers buffer against uncertainty, not against mean travel time. A program that reduces mean PTI from 1.86 to 1.15 but leaves PTI variance unchanged would capture less than one-third of the total reliability benefit. The I2.0 portfolio, by simultaneously addressing managed lane throughput, interchange failure points, and compound closure risk, attacks all three sources of PTI variance and captures the full benefit.
